\uuid{sqCu}
\titre{Ensemble}
\auteur{bodin ; hanani ; mzari}
\organisation{exo7}
\licence{CC BY-NC-SA 4.0}
\datecreate{2020-01-20}
\niveau{L1}
\module{Algèbre}
\chapitre{Logique, ensemble, raisonnement}
\sousChapitre{Ensemble}
\difficulte{2}
\theme{Ensembles, applications, ensemble}
\exo7qcmid{lille-1.ensembles.013}

\contenu{
\begin{question}
\qclassification{100.02,101.01,102.02,102.02}
\qtags{motcle=ensemble}
Soient $A,B$ deux parties d'un ensemble $E$. On note $\overline{A}$ le complémentaire de $A$ dans $E$. Quelle est la bonne réponse ?
\begin{answers}
    \bad{$\overline{A\cup B}=\overline{A}\cup \overline{B}$}
    \good{$\overline{A\cup B}=\overline{A}\cap \overline{B}$}
    \bad{$\overline{A\cup B}=A\cap B$}
    \bad{$\overline{A\cup B}=\overline{A}\cup B$}
\end{answers}
\begin{explanations}
D'abord $x\in A\cup B \Leftrightarrow (x\in A$ ou $x\in B$). Les lois de De Morgan donnent donc que $(x\notin A\cup B)\Leftrightarrow (x\notin A$ et $x\notin B$), c'est-à-dire $\overline{A\cup B}=\overline{A}\cap \overline{B}$.
\end{explanations}
\end{question}
}
